\subsection{可分代数扩张的稳定性性质}

命题8. —— 设E是K的一个扩张，且 \((\mathrm{E}_{i})_{i\in I}\) 是E的一族子扩张，使得 \(\mathrm{E}=\mathrm{K}\left(\bigcup_{i\in I}\mathrm{E}_{i}\right)\) 。若每个扩张 \(E_{i}\) 都是代数且

在K上可分，则E亦然。

这直接由命题6（V，第39页）得出。

命题9. --- 设F是K的一个代数扩张，E是F的一个子扩张。要使 F 在 K 上可分，必要且充分的是 F 在 E 上可分且 E 在 K 上可分。

首先假设F在K上是可分的；那么根据命题1（V，第36页），E在K上也是可分的。此外，F的每个元素在K上都是可分的（V，第39页，命题6），因此在E上也是可分的（V，第39页，推论2），所以F在E上是可分的（V，第39页，命题6）。

反之，假设 F 在 E 上可分，且 E 在 K 上可分。用 x 表示 F 的一个元素，用 \(f \in E[X]\) 表示 x 在 E 上的极小多项式。由于 E 在 K 上是代数的，定理 2（V，第 18 页）表明存在 E 的一个在 K 上有限次的子扩张 \(E'\) ，使得 \(f \in E'[X]\) ；于是 f 同时是 x 在 E 上和在 \(E'\) 上的极小多项式，且由于 x 在 E 上可分（V，第 39 页，命题 6），它在 \(E'(V, p. 39, Prop. 5)\) 上也是可分的。记 \(F' = E'(x)\) ；则 \(F'\) 在 \(E'\) 上是可分的且有限次的，且由于 E 在 K 上可分，\(E'\) 在 K 上可分且有限次（V，第 36 页，命题 1）。因此由 V，第 32 页，推论 2，\(F'\) 在 K 上可分且有限次。因此（V，第 39 页，命题 6）x 在 K 上可分。我们现在已证明 F 的每个元素都在 K 上可分，因此 F 在 K 上可分（V，第 39 页，命题 6）。

命题10. --- 设 E 和 K' 是 K 的同一个扩张的两个子扩张，且设 \(E' = K'(E)\) 。假设 E 在 K 上是代数的，因此 \(E'\) 在 \(K'\) 上是代数的（V，第18页，推论2）。

\begin{enumerate}
\def\labelenumi{\alph{enumi})}
\item
  如果 \(\mathbf{E}\) 在 \(\mathbf{K}\) 上可分，则 \(\mathbf{E}'\) 在 \(\mathbf{K}'\) 上可分。
\item
  反之，如果 \(E'\) 在 \(K'\) 上可分，且 E 与 \(K'\) 在 K 上线性不相交，则 E 在 K 上是可分的。
\end{enumerate}

断言a)直接由命题6（V，第39页）得出。

在假设 \(b\) ) 下，设 \(F\) 为 \(E\) 的一个在 \(K\) 上有限次的子扩张。则 \(F\) 与 \(K'\) 在 \(K\) 上线性不相交，因此 \(K'\) -代数 \(F_{(K')} = K' \otimes_K F\) 同构于 \(K'(F)\) 。由于 \(K'(F)\) 是 \(E'\) 的一个在 \(K'\) 上有限次的子扩张，且 \(E'\) 在 \(K'\) 上是代数的且可分的，\(K'\) -代数 \(K'(F)\) 是平展的。换言之，\(K'\) -代数 \(F_{(K')}\) 是平展的，于是命题4（V，第32页）的推论2表明 \(F\) 在 \(K\) 上是平展的。这样我们就证明了 \(E\) 在 \(K\) 上是可分的。
% semantic-fragments: legacy-input-seam
\relax{}
